The system does not prove a full DQN training run. It does not prove that a
final checkpoint follows from an initialization through all replay batches,
target-network synchronizations, optimizer states, or model-selection decisions.
It also does not prove Adam or full optimizer-state correctness; the tiny update
extension is scoped to simple SGD.

The current backend coverage is still uneven by design. SP1 proof-backed
coverage is limited to the canonical TD MVP vector, Merkle membership including
the generated 1k/10k/100k dataset-size cases, Forward-TD MLP and one-step SGD
tiny canonical vectors, short trace, one-step training update batch size 1,
training fragments with $k\in\{1,4,8\}$, and proof-manifest chunk-chain
aggregation with $T\in\{32,64,128\}$. Distinct minibatch TD and longer
fragments remain semantic-oracle or execute-only evidence unless separately
listed as proof-backed in the provenance tables.

There is no batch-size 4/8/16 training-update proof, no small-network proof
beyond the canonical tiny vectors, and no proof-backed claim for
$k\in\{16,32,128\}$ training fragments. There is also no Atari-scale or
MuJoCo-scale proof and no on-chain verifier.

The benchmarks are small-scale. CartPole and MountainCar are classic-control
tasks, and MountainCar uses a random-policy transition dataset for the second
environment artifact. These choices are appropriate for proving relation
correctness and reproducibility, but they are not evidence of large-scale RL
training performance.

The commitment layer starts after data collection. A verifier can check that a
revealed transition belongs to the committed replay set, but not that the replay
set was collected honestly or that it covers a desired behavior distribution.
Public Minari/D4RL benchmark imports are source-integrity rows only, not honest
public collection proofs.
The Python verifier path also reveals or processes witness data directly and
therefore should not be treated as a privacy-preserving verifier.

Finally, aggregation comes in two forms. Manifest chunk-chain aggregation
covers $T\in\{32,64,128\}$; recursive aggregation, which verifies each child
proof inside the guest, is established for $T\in\{16,32,64\}$ and requires a
CUDA device with roughly $20$~GB of memory. The SP1 proof cost is still
prototype-scale: verification is fast, but proving can take minutes for the
reported relations.
Scaling beyond these scoped relations will require more efficient
neural-network proofs, verified child-proof aggregation, and careful treatment
of optimizer and sampling semantics.
